%% kapa: exam=1 round=20 year=2009 period=1 form=A subject=economics subject_ko=경제원론 seq=31 number=71
%% source: Q-Net 감정평가사 기출문제 (cst00309) artlSeq=5024178 file=extracted/1008/5024178/2009년 제20회 감정평가사 1차시험 기출문제/1교시.hwp
%% conversion: kordoc 4.24.51 -> tools/gen1.py (auto); review=proofread-sonnet; flags=none
\begin{question}{71}
우리나라 국제수지에 관한 설명으로 옳은 것은?
\begin{choices}
  \item 유학생에 대한 해외 송금액 증가는 자본수지 적자 요인이다.
  \item 상품수지와 서비스수지는 동시에 적자를 기록할 수 없다.
  \item 외국인의 우리나라 채권보유 증가는 자본수지 적자 요인이다.
  \item 국내 기업의 해외 건설수주액 증가는 경상수지 적자 요인이다.
  \item 외국인에 대한 주식배당금의 해외 송금은 경상수지 적자 요인이다.
\end{choices}
\end{question}
\answer{5}
